\chapter{Methodology}
\label{ch:methodology}

% Replace this paragraph with an overview of your approach. Cite
% \cite{jones2023} where appropriate.
The approach follows a fixed sequence of steps so that the same run can be
repeated and compared. Each step takes the output of the previous one and
produces a well-defined result, following the structured workflow
pattern~\cite{jones2023}.

% Mermaid system architecture diagram. Replace the flowchart with your own
% system. Keep style= and the caption (wrap a caption containing a comma in
% braces).
\begin{mermaid}[style=technical, width=0.85\linewidth, pos=H, caption={System architecture}]
  flowchart LR
  A[Input data] --> B[Processing]
  B --> C[Model]
  C --> D[Results]
\end{mermaid}

% Replace this paragraph with a sentence linking the diagram to the
% requirements table. Diagrams carry no label, so refer to them by name
% rather than with \ref.
The stages above turn the input data into the reported results, and each
stage is bound to one functional requirement in the table that follows.

% Booktabs requirements table. Replace the rows with your functional
% requirements; keep the caption and the label.
\begin{table}[H]
  \centering
  \caption{Functional requirements}
  \label{tab:requirements}
  \begin{tabular}{@{}llc@{}}
    \toprule
    ID & Requirement & Priority \\
    \midrule
    R1 & Capture the input data & High \\
    R2 & Validate the input & High \\
    R3 & Produce the report & Medium \\
    \bottomrule
  \end{tabular}
\end{table}

Table~\ref{tab:requirements} lists the functional requirements that drive
the design.
